\chapter{Editors: text input internals}
\label{ch:editors}

\begin{aside}
A text input feels simple: type, see characters. The
implementation hides a cursor, a selection, an undo stack,
clipboard interaction, history, IME quirks, and Unicode word
boundaries. Each is a small thing; together they are an
editor.
\end{aside}

\section{The buffer}

The simplest buffer is a string:

\begin{lstlisting}
std::string text;
size_t cursor;
\end{lstlisting}

Insertion at cursor: \code{text.insert(cursor, c); ++cursor}.
Deletion: \code{text.erase(cursor, 1)}.

For long buffers, a gap buffer or a rope is faster (O(1)
insertion at gap or splay point). Single-line inputs use a
string; multi-line editors use a piece table or rope.

\section{Cursor and selection}

The selection is a range from anchor to cursor.
When no selection, anchor equals cursor. Movement
collapses the selection unless Shift is held.

\section{Word boundaries}

Ctrl-Left jumps to the previous word boundary.
Boundaries follow Unicode TR\#29 word-segmentation rules:
between letters and punctuation, between digits and
letters, etc.

\maya{} implements a small state machine.

\section{Undo}

Every mutation pushes onto an undo stack:

\begin{lstlisting}
struct Edit {
    size_t pos;
    std::string old_text;
    std::string new_text;
};
std::vector<Edit> undo_stack;
\end{lstlisting}

Undo pops, reverses the edit. Redo pushes the reverse onto
a redo stack.

\section{Grouping}

Each keystroke as a separate undo entry is annoying. \maya{}
groups consecutive insertions of word characters into one
undoable unit. A boundary (space, punctuation, cursor move)
starts a new group.

\section{Clipboard}

OSC-52 lets terminals access the system clipboard:

\begin{lstlisting}
ESC ] 52 ; c ; BASE64 BEL
\end{lstlisting}

Letter c for clipboard, p for primary selection. The
terminal copies the base64-decoded text.

Pasting comes via bracketed paste, not via OSC-52.

\section{IME}

Input methods for CJK languages compose characters before
commit. The terminal handles the popup; \maya{} sees the
committed string. No special handling required.

\section{History}

A separate history (up/down arrow recalls previous values)
is common for command lines. The history stack persists
across sessions in a file.

\section{Multi-line}

A multi-line editor needs:
\begin{itemize}[leftmargin=*]
  \item Per-line wrapping.
  \item Vertical cursor movement (preserve column).
  \item Scrolling.
  \item Line-numbered display (optional).
\end{itemize}

\maya{} provides \code{multiline\_input} that wraps the
single-line input with these affordances.

\section{Recap}

\begin{recap}
\begin{itemize}[leftmargin=*]
  \item String buffer for short; rope for long.
  \item Selection = anchor + cursor.
  \item Word boundaries via TR\#29.
  \item Undo grouped by word.
  \item OSC-52 for clipboard.
  \item IME is the terminal's problem.
\end{itemize}
\end{recap}

\section*{Workshop}

\begin{enumerate}[leftmargin=*]
  \item Build a text input. Add Ctrl-Left/Right word jumps.
  \item Implement undo/redo with grouping.
  \item Wire OSC-52 copy on Ctrl-C.
\end{enumerate}
